\section*{अध्याय \ref{ch:g2:from-seed-to-plant} --- बीज से पौधे तक}

\begin{solution}{exo:g2:from-seed-to-plant:1}
एक नन्हा \omterm{def:g1:plants-around-us:plant}{पौधा} जो शुरू हो चुका है (\omterm{def:g1:plants-around-us:plant}{पौधे} का \omterm{def:g2:from-seed-to-plant:seed}{भ्रूण}, जिसमें एक \omterm{def:g1:plants-around-us:parts}{जड़} और एक
\omterm{prop:g2:from-seed-to-plant:cycle}{अंकुर} की शुरुआत होती है) और उसे खिलाने के लिए भोजन का एक भंडार।
\end{solution}

\begin{solution}{exo:g2:from-seed-to-plant:2}
\omterm{def:g2:from-seed-to-plant:germination}{अंकुरित होना} यानी जागकर बढ़ना शुरू करना: \omterm{def:g2:from-seed-to-plant:seed}{बीज} फूलता है, उसका छिलका
फटता है, और वह \omterm{def:g1:plants-around-us:plant}{पौधा} बनना शुरू कर देता है। सबसे पहले \omterm{def:g1:plants-around-us:parts}{जड़} बाहर आती है,
और नीचे की ओर गोता लगाती है।
\end{solution}

\begin{solution}{exo:g2:from-seed-to-plant:3}
पानी और गरमाहट। प्रकाश सूची में नहीं है --- \omterm{def:g2:from-seed-to-plant:seed}{बीज} मिट्टी के नीचे,
अँधेरे में भी मज़े से \omterm{def:g2:from-seed-to-plant:germination}{अंकुरित होता है}।
\end{solution}

\begin{solution}{exo:g2:from-seed-to-plant:4}
वह \omterm{def:g2:from-seed-to-plant:seed}{बीज} के भीतर भरे भोजन-भंडार पर जीता है --- सेम के वही दो मोटे
टुकड़े। यह भंडार \omterm{prop:g2:from-seed-to-plant:cycle}{अंकुर} को तब तक खिलाता है जब तक उसकी हरी \omterm{def:g1:plants-around-us:parts}{पत्तियाँ}
प्रकाश तक न पहुँच जाएँ और वह ख़ुद न खा सके।
\end{solution}

\begin{solution}{exo:g2:from-seed-to-plant:5}
जैसे: सेम, मटर, मसूर, चावल, सूरजमुखी के \omterm{def:g2:from-seed-to-plant:seed}{बीज}, सेब के दाने, आड़ू की
गुठली।
\end{solution}

\begin{solution}{exo:g2:from-seed-to-plant:6}
क्योंकि \omterm{def:g2:from-seed-to-plant:seed}{बीज} को पानी के साथ-साथ गरमाहट भी चाहिए: सर्दियों की ठंडी
मिट्टी में वह \omterm{def:g2:from-seed-to-plant:germination}{अंकुरित होने} के बजाय इंतज़ार करता, या सड़ जाता। वसंत
वही गरम मिट्टी लाता है जिसका \omterm{def:g2:from-seed-to-plant:seed}{बीज} इंतज़ार कर रहा है।
\end{solution}

\begin{solution}{exo:g2:from-seed-to-plant:7}
\omterm{def:g2:from-seed-to-plant:seed}{बीज}, \omterm{prop:g2:from-seed-to-plant:cycle}{अंकुर}, पत्तेदार \omterm{def:g1:plants-around-us:plant}{पौधा}, \omterm{def:g1:plants-around-us:parts}{फूल} वाला \omterm{def:g1:plants-around-us:plant}{पौधा}, नए सेमों वाली फली।
\end{solution}

\begin{solution}{exo:g2:from-seed-to-plant:8}
दोनों की \omterm{def:g1:plants-around-us:parts}{जड़ें} नीचे की ओर बढ़ती हैं और दोनों के \omterm{prop:g2:from-seed-to-plant:cycle}{अंकुर} ऊपर की ओर, चाहे
सेम जैसे भी पड़े हों। \omterm{prop:g2:from-seed-to-plant:cycle}{अंकुर} हमेशा ऊपर और नीचे पहचान लेता है --- ठीक
वही जो टिप्पणी में बताया गया है।
\end{solution}

\begin{solution}{exo:g2:from-seed-to-plant:9}
दराज़ में \omterm{def:g2:from-seed-to-plant:seed}{बीजों} को पानी नहीं मिला, इसलिए वे अंकुरित नहीं हो सके; पर
अच्छा \omterm{def:g2:from-seed-to-plant:seed}{बीज} पानी के बिना मरता नहीं --- वह इंतज़ार करता है। साल भर बाद
पानी और गरमाहट मिलते ही इंतज़ार करते \omterm{def:g2:from-seed-to-plant:seed}{बीज} सामान्य रूप से जाग उठे।
\end{solution}

\begin{solution}{exo:g2:from-seed-to-plant:10}
\omterm{def:g2:from-seed-to-plant:seed}{बीज} का भोजन-भंडार --- सेम के दो मोटे टुकड़े --- ज़र्दी की भूमिका
निभाता है। \omterm{def:g2:animals-grow-up:born}{अंडा} और \omterm{def:g2:from-seed-to-plant:seed}{बीज}, दोनों किसी नए \omterm{def:g1:living-or-not:living}{सजीव} की शुरुआत हैं, और दोनों
में वह भोजन भरा होता है जो नन्हे जीव को ख़ुद खा सकने से पहले चाहिए।
\end{solution}
